% TOP_PROBLEM: {"id": 1100404, "title": "Questions in Geometric Group Theory — Q 4.4", "category_id": 17, "category": {"id": 17, "name": "group_theory", "display_name": "Group Theory", "description": "Problems about groups, group actions, representations, and related algebraic structures.", "slug": "group-theory", "order_index": 17, "created_at": "2026-07-31T15:26:25.671Z"}, "status": "open", "problem_number": "AMR-010-0404", "set_id": 13}
% FIRST_POSTED: 2026-10-02
% ENTRY_KIND: statement_only
% UnsolvedMath ID 1100404; source catalogue code AMR-010-0404
% !TeX program = lualatex
% Definition and sources only; this file is not an LLM solution attempt.
\documentclass[11pt,a4paper]{article}
\usepackage[margin=25mm]{geometry}
\usepackage{fontspec}
\setmainfont{lmroman10-regular.otf}[BoldFont=lmroman10-bold.otf,
 ItalicFont=lmroman10-italic.otf,BoldItalicFont=lmroman10-bolditalic.otf]
\usepackage{amsmath,amssymb,mathtools}
\usepackage{unicode-math}
\setmathfont{latinmodern-math.otf}
\AtBeginDocument{\let\setminus\smallsetminus}
\usepackage{xurl,hyperref}
\hypersetup{colorlinks=true,urlcolor=blue}
\setcounter{secnumdepth}{0}
\setlength{\parindent}{0pt}
\setlength{\parskip}{5pt}
\setlength{\emergencystretch}{3em}
\begin{document}
\section*{Questions in Geometric Group Theory — Q 4.4}\label{1100404}
{\small UnsolvedMath ID 1100404 · AMR-010-0404 · Group Theory · AMR Open Problem Lists}
\subsection{Definitions and mathematical statement}
Let $G$ be a finitely presented one-ended group. One-ended means that deleting any finite vertex set from a Cayley graph leaves exactly one infinite connected component. Suppose $G$ acts by homeomorphisms on a compact connected metric space $X$ containing at least three points. Put
\[
 X^{(3)}=\{(x_1,x_2,x_3)\in X^3:x_i\ne x_j\text{ for }i\ne j\}.
\]
The action is a convergence action if it acts properly on this triple space: for compact $K\subset X^{(3)}$, only finitely many $g\in G$ satisfy $gK\cap K\ne\varnothing$. A point $p\in X$ is a cut point when $X\setminus\{p\}$ is disconnected. Local connectedness means every point has a basis of connected open neighborhoods.

Determine the consequences for $G$ and its action when $X$ has a cut point. In particular, determine under what hypotheses this forces a nontrivial splitting of $G$, and whether $X$ must be locally connected in the setting of the source question.

The topology concerns the given convergence space, not an arbitrarily chosen Cayley-graph compactification. The motivating hyperbolic-boundary theorem assumes that $X$ is the Gromov boundary of a hyperbolic group; that conclusion is not silently applied to a general convergence action. If a minimality or geometric-finiteness assumption is necessary, identifying it is part of resolving the stated structural problem.
\subsection{Short English statement}
What do cut points and local connectedness of a convergence-action space tell us about a one-ended finitely presented group?
\subsection{Sources}
\begin{itemize}
\item[S1] ULAM.AI and the UnsolvedMath Contributors, supplied problems.json, record 1100404 (AMR-010-0404), AMR Open Problem Lists. Catalogue identity and source transcription; CC BY 4.0.
\url{https://huggingface.co/datasets/ulamai/UnsolvedMath}.
\item[S2] Bestvina - Questions in Geometric Group Theory (pdf) (2004), Question 4.4 (PDF page 12). Original source indexed by AMR-010-0404.
\url{https://www.math.utah.edu/~bestvina/eprints/questions-updated.pdf}.
\end{itemize}
\end{document}
